\chapter{System Testing and Evaluation}

\section{Testing Objectives}

The main objective of testing the ENDOTECH system is to verify the correctness, reliability, and efficiency of the implemented functionality.

\section{Testing Methods}

The following testing methods were applied:

\begin{itemize}
    \item Functional testing;
    \item Integration testing;
    \item Manual testing;
    \item Usability testing.
\end{itemize}

\section{Functional Testing}

The following functionalities were tested:

\begin{itemize}
    \item User authentication using Firebase Authentication;
    \item Inventory management operations (add, update, delete);
    \item Sales processing and stock updates;
    \item Real-time data synchronization using Firestore.
\end{itemize}

All features were tested successfully and performed as expected.

\section{Performance Testing}

Performance testing showed:

\begin{itemize}
    \item Fast response time for database queries;
    \item Stable operation under normal load;
    \item Efficient real-time updates using Firebase.
\end{itemize}

\section{Usability Testing}

Usability testing results:

\begin{itemize}
    \item The interface is intuitive and easy to use;
    \item Navigation is clear and efficient;
    \item Users can perform tasks quickly.
\end{itemize}

\section{Limitations}

Some limitations were identified:

\begin{itemize}
    \item Limited number of test devices;
    \item Small dataset size;
    \item Limited number of users in usability testing.
\end{itemize}

\section{Conclusion}

The testing results confirm that the ENDOTECH system is stable, efficient, and suitable for real-world use.